\section*{Scriptural Date and Location}

{\footnotesize Seven distinct passages are directly appointed, one for each
scriptural element; the three orations quote no Scripture. Psalms carry the
missal's Vulgate number, the common modern number in parentheses.\par}

% Full-width passage dossiers keep every generated relation and alternative
% readable without compressing the Date column or reducing the text size.
\begingroup
\footnotesize
\renewcommand{\arraystretch}{1.02}
\setlength{\LTpre}{0pt}
\setlength{\LTpost}{0pt}
\begin{longtable}{@{}>{\raggedright\arraybackslash}p{\linewidth}@{}}
\toprule
\textbf{Offertory --- Ps. 39:2, 3 \latin{et} 4 (Eng. 40:1--3)}. Title \latin{In finem. Psalmus ipsi David}; no occasion named\\*
\chronodate{offertory}{\chronologyannotation{offertory}}\\*
The unsung title at Vulgate v.~1 names David without an occasion, place or life stage. Corbett's ``usual chronology'' (\work{Catholic Encyclopedia}, IV, p.~642) supplies a regnal reference, not the psalm's date, alongside a later reckoning. The Biblical Commission calls the titles ancient and David the Psalter's chief author (\work{AAS} 2 (1910) 354--355). Drum (XII, ``Psalms'') dates none of this Mass's four psalms; the critical bound applies to the whole Psalter, whose individual psalms cannot be dated securely.\\
\midrule
\textbf{Introit --- Ps. 85:1, 2--3; v.~4 (Eng. 86:1, 2--3, 4)}. Title \latin{Oratio ipsi David}; the one Davidic psalm in Book~III\\*
\chronodate{introit}{\chronologyannotation{introit}}\\*
The antiphon omits the ascription, the Clementine's only Davidic title in Book~III. Neither title nor psalm names an occasion, place or life stage. The regnal reference, 1910 responsa and whole-Psalter bound are those of the Offertory.\\
\midrule
\textbf{Gradual --- Ps. 91:2--3 (Eng. 92:1--2)}. Title \latin{Psalmus cantici, in die sabbati}; no place, no occasion\\*
\chronodate{gradual}{\chronologyannotation{gradual}}\\*
The unsung title at Vulgate v.~1 names the sabbath but no author. Composition is bounded only as the whole Psalter's is.\\
\midrule
\textbf{Alleluia --- Ps. 94:3 (Eng. 95:3)}. Title \latin{Laus cantici ipsi David}; summons vv.~1--7, warning vv.~8--11\\*
\chronodate{alleluia}{\chronologyannotation{alleluia}}\\*
The verse begins the summons' reasons; the warning at v.~8 is unsung. The title names David without an occasion. It is part of Vulgate v.~1, so English verse numbers agree. The regnal reference and Psalter bound are those of the Offertory.\\
\midrule
\textbf{Gospel --- Lk. 7:11--16}. Naim, in Galilee; the Third Gospel written by Luke for Gentile converts\\*
\chronodate{gospel}{\chronologyannotation{gospel}}\\*
\textit{Event: at the gate of the city called Naim, next in order after the
healing of the centurion's servant at Capharnaum (Lk.\ 7:1--10), with the
disciples and a great multitude walking with Christ and the town's own crowd
walking with the widow.}\\*
Maas (\work{Catholic Encyclopedia}, VIII, ``Jesus Christ'') dates the event; X, ``Naim'', locates it south of Thabor near Endor in Eusebius and Jerome's \work{Onomasticon}, and identifies Nain on Jebel Dahy's north-western ridge. Aherne (IX, ``Gospel of Saint Luke'') supplies authorship, Gentile audience and purpose. Composition alternatives follow Durand (XIV, ``The New Testament'') and the 1912 Biblical Commission (\work{AAS} 4, 463--465). No independent modern critical composition date is established here.\\
\midrule
\textbf{Communion --- Io. 6:52 (6:51 in the modern series)}. The synagogue at Capharnaum (Jn.\ 6:60)\\*
\chronodate{communion}{\chronologyannotation{communion}}\\*
\textit{Event: the Bread of Life discourse, spoken in the synagogue at
Capharnaum on the day after the multiplication of the loaves and near the Pasch
(Jn.\ 6:4, 6:22--24, 6:60).}\\*
Maas dates the event; Fonck (\work{Catholic Encyclopedia}, VIII, ``Gospel of St John'') treats authorship, audience and writing place. Fonck and Durand (XIV, ``The New Testament'') supply the disputed composition dates; no modern critical date is established here. Clementine, Douay and Missal number the clause 52; the Greek Fathers and translations from them number it 51.\\
\midrule
\textbf{Epistle --- Gal. 5:25--26; 6:1--10}. Galatia's churches (Gal.\ 1:2); writing place undetermined\\*
\chronodate{epistle}{\chronologyannotation{epistle}}\\*
Paul names himself and Galatia's churches (1:1--2). A~Lapide's preferred date comes from the prooemium and \latin{Argumentum} of \work{Commentaria in omnes divi Pauli epistolas} (Antwerp, 1614). Aherne (\work{Catholic Encyclopedia}, VI, ``Galatians'') and Prat (XI, ``St Paul'') supply four alternatives. The nine-year spread is traditional disagreement; no modern critical date is established here.\\
\bottomrule
\end{longtable}
\endgroup

\clearpage
